\documentclass{svproc}
%
\usepackage{graphicx}%
\usepackage{multirow}%
\usepackage{amsmath,amssymb,amsfonts}%
\usepackage{mathrsfs}%
\usepackage[title]{appendix}%
\usepackage{xcolor}%
\usepackage{textcomp}%
\usepackage{manyfoot}%
\usepackage{booktabs}%
\usepackage{algorithm}%
\usepackage{algorithmicx}%
\usepackage{algpseudocode}%
\usepackage{listings}%
\usepackage{url}
\def\UrlFont{\rmfamily}

\renewcommand{\topfraction}{.99}
\renewcommand{\dbltopfraction}{.99}
\renewcommand{\bottomfraction}{.99}
\renewcommand{\textfraction}{.01}
\renewcommand{\floatpagefraction}{0.9}
\renewcommand{\dblfloatpagefraction}{0.9}

\usepackage[normalem]{ulem}
\usepackage{bm}
\usepackage{xspace}
\usepackage{subcaption}

\usepackage{color}
\def\green#1{\textcolor[rgb]{0,0.6,0}{#1}}
\def\black#1{\textcolor[rgb]{0,0,0}{#1}}
\def\dgreen#1{\textcolor[rgb]{0.02,0.50,0.05}{#1}}
\def\gray#1{\textcolor[rgb]{0.5,0.5,0.7}{#1}}
\def\red#1{\textcolor{red}{#1}}
\def\blue#1{\textcolor{blue}{#1}}
\def\Blue{\color[rgb]{0,0,1}}
\def\Red{\color[rgb]{1,0,0}}
\def\Dgreen{\color[rgb]{0.02,0.55,0.04}}
\def\Black{\color[rgb]{0,0,0}}
\def\ignore#1{}

\usepackage[normalem]{ulem}
\usepackage{bm}
\usepackage{booktabs}
\usepackage{xspace}
\usepackage{subcaption}

% track-changes helpers (kept from last year's setup)
\newcommand{\add}[1]{\textcolor{blue}{#1}}
\newcommand{\erase}[1]{\textcolor{red}{\sout{#1}}}
\def\subst#1#2{\add{#1}{ \erase{#2}}}

% parameters definition (same convention as docs/journal.tex)
\def\post{\rho}                 % posting probability
\def\BoundedConfidence{\epsilon} % bounded confidence (rejection) threshold
\def\feed{F}                    % feed
\def\Stubbornness{\mu}           % stubbornness
\def\maxfollower{S_{\it fol}}    % max followee capacity
\def\maxfeed{S_{\it fd}}         % feed capacity
\def\VCR{\kappa}                 % vocal comfort radius (expression threshold)

\begin{document}
\mainmatter

\title{When Do Moderate Hubs Fuel Extreme Opinion Activation? \\ A Two-Layer Homophily Mechanism in Adaptive Social Networks}
\titlerunning{Moderate Hubs and Two-Layer Homophily}

\author{Tomoya Takeda\inst{1}, Fujio Toriumi\inst{2}, Toshiharu Sugawara\inst{1}}
\authorrunning{Tomoya Takeda, Fujio Toriumi and Toshiharu Sugawara} % abbreviated author list (for running head)

\institute{Department of Computer Science and Communications Engineering,
\\Waseda University, 3-4-1 Okubo, Shinjyuku-ku, 169-8555, Tokyo, Japan.\\
\email{tomoya.leo9@fuji.waseda.jp}, \email{sugawara@waseda.jp}
\and
Department of Systems Innovation, The University of Tokyo, \\7-3-1
Hongo, Bunkyo-ku, 113-8656, Tokyo, Japan.\\
\email{tori@sys.t.u-tokyo.ac.jp}
}

\maketitle

\begin{abstract}\Dgreen
With the rapid expansion of web communication platforms, hub users
have become critical actors in online opinion dynamics, yet their
effects on polarized posting behavior remain poorly understood. Using
an agent-based model on heavy-tailed adaptive social networks, we
study how moderately positioned hubs shape opinion-sharing
behavior. Active moderate hubs consistently boost moderate-camp
posting while suppressing near-extreme posting. Their influence on the
ideologically distant extreme camp, however, is conditional: moderate
hubs can act as mediators or facilitators of activation depending on
network conditions. We trace this regime shift to a two-layer
homophily mechanism—structural and perceptual—which determines
whether the two camps are perceived as one group (mediation) or as
separate communities (facilitation). Moderate hubs are therefore not
universally stabilizing bridges; their brokerage role flips between
mediation and facilitation depending on how homophilic the two camps
are perceived to be.
\ignore{With the rapid expansion of web communication platforms, hub users have become critical actors in online opinion expression dynamics; however, their effects on polarized posting behavior remain poorly understood. Using an agent-based model on heavy-tailed adaptive social networks, we study how moderately positioned hubs shape opinion-sharing behavior. Across a comprehensive suite of experiments, we observe a consistent pattern in which active moderate hubs boost moderate-camp posting while suppressing near-extreme posting. Their influence on the ideologically distant extreme camp, however, is highly complex: moderate hubs can act as mediators or facilitators of activation depending on network conditions. We trace this regime shift to a two-layer homophily mechanism, structural and perceptual, which determines whether the two camps are perceived as one group (mediation) or as separate communities (facilitation). Moderate hubs are therefore not universally stabilizing bridges; their brokerage role flips between mediation and facilitation depending on how homophilic the two camps are perceived to be.}
\Black
\keywords{Social networks, Opinion dynamics, Agent-based simulation,
  Social media, Computational Social Science} 
\end{abstract}

\Blue


\section{Introduction}\label{sec:intro}
Social media and web communication platforms have become the primary
venues for various styles of speech and discussion, rather than merely
for sharing individual ideas, thereby exerting a significant effect on
public opinion formation. Unlike traditional face-to-face
communication, hub users or influencers --- followed by a large number
of users --- can create posts that are disproportionately amplified
across a platform's population, giving them immense influence on a
given opinion environment. The role and effect of these hubs have been
well studied through agent-based simulation and data mining,
particularly from the perspective of information diffusion. However,
their effects on polarized, controversial opinion dynamics are not yet
fully understood, specifically regarding moderately positioned hubs as
opposed to strongly identified opinion leaders.
\par


Therefore, we investigate how the sharing activity of these moderate
hubs affects polarized opinion dynamics. Using an agent-based
simulation on heavy-tailed adaptive social networks, we find that
active moderate hubs do not uniformly mediate extreme-camp posting:
they consistently boost posting activity among aligned moderate users
while suppressing that of the ideologically nearby near-extreme camp,
but their effect on the ideologically distant extreme camp is
conditional. The results show that this conditional effect manifests
as a clear regime shift governed by perceptual and structural
parameters.
\par


Although mechanistically distinct, the effects of these two parameters
appear to converge on a single, unifying homophily logic: when the
near-extreme and extreme camps are effectively perceived as one
homogeneous group, suppressing near-extreme posting also mediates
extreme-camp activity. When they are instead perceived as separate,
distinct communities, moderate hubs' suppression of near-extreme
posting does not reach the extreme camp. Instead, it isolates the
extreme camp into a more pure, self-contained echo chamber, which
facilitates, rather than dampens, its activation. This perception of
homophily versus heterophily is governed by two distinct layers: a
perceptual layer, defined by how wide an opinion difference users find
comfortable to engage with, which determines whether reached content
registers as similar enough to elicit participation; and a structural
layer, defined by how readily agents disconnect from ties they no
longer tolerate, which shifts the threshold of opinion difference that
must be tolerated for the two camps to be treated as a single group.
\par

\Black

%[How these two layers relate to one another]

The role of bridge nodes, or brokers, has long been studied, yet its
effect on social media remains highly complex due to the varied
dynamics such platforms exhibit.This study offers a novel, mechanistic
account of the conditional role of brokerage nodes.

\Blue

\section{Related Work}\label{sec:related}
Although social networks have long been studied as static structures,
most notably in seminal work on small-world topology~\cite{watts1998},
they have recently been modeled as systems whose topology coevolves
with the opinions and behaviors of the individuals embedded in
them. Building on the classic \emph{voter model}~\cite{clifford1973},
Vazquez et al.~\cite{vazquez2008} used an {\em agent-based model}
(ABM) to analyze coevolutionary dynamics based on two coupled rules
--- state updating and edge rewiring --- revealing a phase transition
between connected and fragmented networks governed by degree
settings. The study of opinion dynamics is one of the major fields
where this style of coevolving network model has been fruitfully
applied: Kozma and Barrat~\cite{kozma2008} modeled how agents
adaptively rewire their social ties in response to disagreement,
showing the dominant effect of homophilic rewiring on the resulting
network structure. This coevolving network model is based on {\em bounded
confidence}, formalized by Deffuant et al.~\cite{deffuant2000} and
Hegselmann and Krause~\cite{hegselmann2002}, in which an agent updates
its opinion only by averaging with others whose opinion falls within a
fixed tolerance window. With the rapid proliferation of social media
in recent years, public discourse has shifted significantly from
traditional face-to-face communication. In line with this
homophilic-rewiring mechanism, Sasahara et al.~\cite{sasahara2020}
used an ABM to confirm that social influence combined with the
unfollowing process successfully reproduces the politically polarized
follow network observed in real-world social media
data~\cite{conover2011}.
\par


Whether such bridging roles are inherently stabilizing, however, is
far from settled in the context of social media. Classic network
theory established that ties bridging otherwise disconnected
communities act as a source of connectivity and innovation rather than
redundancy, whether characterized as \emph{weak ties} diffusing
influence across group boundaries~\cite{granovetter1973} or as
\emph{broker} nodes spanning structural holes~\cite{burt1992}. This
bridging role was later quantified via \emph{betweenness centrality},
defined as the proportion of the shortest paths between node pairs
passing through a given node~\cite{freeman1977}. On social media,
however, this integrative reading does not hold unconditionally: Cha
et al.~\cite{cha2010} observed that highly followed users do not
necessarily exert a larger influence, which is an important caveat for
influence maximization in social media marketing --- and subsequent
ABM studies found a hub's effect on the surrounding population to be
conditional rather than uniformly integrative. Using an ABM, Helfmann
et al.~\cite{helfm2023} showed that hub users who adopt extreme
attitudes can gain more followers on social media. Also with an ABM,
Betts et al.~\cite{betts2023} found that, unlike an extremist
influencer whose activity always accelerates polarization, a neutral
influencer's effect on polarization inverts depending on its own
tolerance for disagreeing views; that is, brokerage by a moderate hub
is not inherently stabilizing.
\par


As reviewed above, most of these studies focus on opinion positions
themselves as the primary outcome, whether analyzing an influencer's
own attitude~\cite{helfm2023} or a population's aggregate
polarization~\cite{betts2023,sasahara2020}. On social media, however,
influence operates not only through changes in private opinion but
also through the decision to express or withhold it. Following the
\emph{spiral-of-silence} tradition, in which individuals modulate
whether they voice an opinion according to their perceived opinion
climate~\cite{noelle-neumann1974}, Krassa~\cite{krassa1988} used an
ABM to test this dynamic through social groups and selective
perception. Yet, even simulation-based studies focusing on social
media tend to treat opinion position as the primary outcome, leaving
expression secondary. Therefore, we take posting, or vocality, as our
primary outcome throughout, investigating under what conditions
moderate hubs' ostensibly mediating role produces complex effects on
opinion expression across groups.
\par


\section{Model}\label{sec:model}
\subsection{Agent-Based Model}\label{sec:model-abm}

We simulate opinion and posting dynamics on a directed graph
$G=(V,E)$, where $V=\{1,2,\dots,N\}$ is the set of $N$ agents
representing individual social media users and $E\subseteq V\times V$
is the set of directed follow edges, so that $(i,j)\in E$ means agent
$i$ follows agent $j$, with a corresponding adjacency matrix
$\mathbf{A}$ whose $(i, j)$-th entry is $\mathbf{A}_{ij}=1$ if
$(i,j)\in E$ and $\mathbf{A}_{ij}=0$ otherwise. This follow graph
evolves through the 
social media behaviors introduced in
Section~\ref{sec:model-behaviors}. The dynamics proceed in discrete
time steps $t\in\{0,1,2,\dots\}$.
\par

Agent $i\in V$ holds an opinion $O_i(t)\in[-1,1]$, representing
its stance on a single controversial issue at time $t$; the boundaries
$-1$ and $1$ denote the two most extreme opposing positions, and $0$
denotes a neutral stance. Hereafter, we group agents into three camps
according to their opinion strength by dividing $[-1,1]$ into five
equal-width bins: agents with $\vert O_i(t)\vert <0.2$, in the
innermost bin, form the \emph{moderate} camp; agents with
$0.2\le\vert O_i(t)\vert<0.6$, in either of the two intermediate
bins, form the \emph{near-extreme} camp; and agents with
$\vert O_i(t)\vert\ge0.6$, in either of the two outer bins, form
the \emph{extreme} camp, with their ideological alignment determined
by the sign of $O_i(t)$.
% As is standard in the opinion dynamics literature,
% we compress an agent's stance to this single dimension, 
% under the assumption that a specific controversial 
% topic can be usefully represented along one axis 
% even though real-world opinions are inherently 
% multidimensional~\cite{deffuant2000,hegselmann2002}.
Agent's initial opinion $O_i(0)$, which
we call its \emph{intrinsic opinion} representing a latent attitude,
is drawn independently from a Gaussian distribution, resampling any
draw that falls outside $[-1,1]$.
\par


Beyond its opinion, each agent is characterized by two further
internal states, summarized in Table~\ref{tab:state}: a posting
probability $\post_i(t)$, governing how actively it posts, and a
{\em bounded confidence threshold} $\BoundedConfidence_i(t)$, governing
which content from other agents it tolerates well enough to remain
connected with.
\par

% Every agent additionally carries a fixed
% stubbornness $\Stubbornness$, representing the extent
% to which it adheres to its own
% intrinsic opinion $O_i(0)$ rather than
% being moved by the opinions it encounters;
% in the experiments reported here $\Stubbornness$ is
% held at a constant value across the population,
% though the implementation supports drawing it
% heterogeneously per agent.

\begin{table}[t]\Blue
\centering
\caption{Agent state variables and their value
  ranges.}\label{tab:state}
\tabcolsep=5pt
\begin{tabular}{llll}
\toprule
State variable & Symbol & Range & Initial value \\
\midrule
Opinion & $O_i(t)$ & $[-1.0,\ 1.0]$ & $\mathcal{N}(0.0,\ 0.6^2)$ \\
Posting probability & $\post_i(t)$ & $[0.01,\ 1.0]$ & $0.26$ \\
Bounded confidence threshold & $\BoundedConfidence_i(t)$ & $[0.2,\ 1.0]$ & $0.8$ \\
\bottomrule
\end{tabular}\Black
\end{table}



\subsection{Modeling Social Adaptation}\label{sec:model-adaptation}
Each agent updates its internal state only when it accesses the
platform, an event that occurs independently at each time step with
access probability $P_a$. Therefore, every rule introduced below is
understood to fire once per access rather than once per step.
\par


\paragraph{Social influence.} 
Upon accessing the platform, agent $i$ reads all posts in its current
feed $\feed_i(t)$, which is defined as the set of posts delivered to
$i$ during that access. It then updates its opinion by averaging the
subset of posts within its tolerance threshold. This average is
blended with the agent's intrinsic opinion $O_i(0)$ according to its
stubbornness parameter $\Stubbornness\in[0,1]$, which acts as a fixed
weight on its intrinsic opinion relative to social influence, as
expressed in Eq.~\eqref{eq:social-influence}.
%
\begin{equation}
O_i(t) = \Stubbornness\cdot O_i(0) + (1-\Stubbornness)\cdot\frac{\sum_{p_j\in \feed_i(t):\,|O(p_j)-O_i(t)|<\BoundedConfidence_i(t)} O(p_j)}{\big|\{p_j\in \feed_i(t):\,|O(p_j)-O_i(t)|<\BoundedConfidence_i(t)\}\big|}.
\label{eq:social-influence}
\end{equation}
%
If no post in the feed
falls within $\BoundedConfidence_i(t)$, the agent's opinion remains
unchanged.
% for that access.



\paragraph{Confidence decline and recovery.} 
Bounded confidence $\BoundedConfidence_i(t)$ is adjusted at every
access toward a target level determined by how much of the feed the
agent tolerates, with out-of-tolerance posts narrowing this
threshold. Writing $r_i(t)$ for the fraction of $i$'s feed that falls
within its current threshold, the target level is defined as
%
\begin{equation} 
  \BoundedConfidence_i^{\ast}(t) =
  \BoundedConfidence_{\mathrm{floor}} + (\BoundedConfidence_{\mathrm{ceil}}-\BoundedConfidence_{\mathrm{floor}})\cdot r_i(t), 
\end{equation} 
%
and the threshold is relaxed toward it according to 
%
\begin{equation} 
  \BoundedConfidence_i(t{+}1) = \BoundedConfidence_i(t) +
  \begin{cases}
    \delta_{\BoundedConfidence},
    \big(\BoundedConfidence_i^{\ast}(t)-\BoundedConfidence_i(t)\big)
    & \text{if } \BoundedConfidence_i^{\ast}(t) <
    \BoundedConfidence_i(t), \\[2pt]
    \gamma_{\BoundedConfidence},
    \big(\BoundedConfidence_i^{\ast}(t)-\BoundedConfidence_i(t)\big)
      & \text{otherwise},
  \end{cases}\label{eq:bc} 
\end{equation} 
%
clipped to
$[\BoundedConfidence_{\mathrm{floor}},\BoundedConfidence_{\mathrm{ceil}}]$,
where $\delta_{\BoundedConfidence}$ and $\gamma_{\BoundedConfidence}$
($\delta_{\BoundedConfidence}\gg\gamma_{\BoundedConfidence}$) 
reflects the well-documented asymmetry by which negative experiences
carry greater psychological weight~\cite{baumeister2001}.


\paragraph{Social reinforcement and vocality.} 
Whether an agent tolerates a post in the bounded confidence sense
described above is conceptually distinct from whether that content
emboldens the agent to express its opinion. Therefore, a distinct
second criterion is introduced. A {\em post} $p_j$ by agent $j$ is
defined as {\em comfortable} for agent $i$ if $|O_i(t)-O(p_j)|<\VCR$, where the vocal comfort
radius $\VCR$ is a population-wide threshold, deliberately set apart
from the agent-specific threshold $\BoundedConfidence_i(t)$. The two
criteria capture different functional mechanisms:
$\BoundedConfidence_i(t)$ acts as a rejection threshold, governing how
much disagreement an agent tolerates before disengaging, whereas
$\VCR$ serves as an expression threshold that determines how much
agreement with the surrounding opinion climate is required before
posting.
% An agent can in 
% principle be narrowly tolerant of disagreement 
% yet still easily emboldened to speak, or the reverse. 
Letting $c_i(t)$ denote the fraction of $i$'s feed that is
comfortable, the posting probability relaxes toward a target value
determined linearly by $c_i(t)$:
%
\begin{equation}
\post_i^{\ast}(t) = \post_{\min} + (\post_{\max}-\post_{\min})\cdot c_i(t), \qquad
\post_i(t{+}1) = \post_i(t) + \eta\,\big(\post_i^{\ast}(t)-\post_i(t)\big),
\end{equation}
%
where the updated value is clipped to
$[\post_{\min},\post_{\max}]$. As with the confidence-decline rule
above, this graded formulation ensures that every comfort level yields
a distinct, stable equilibrium vocality.
\par


\subsection{Modeling Behaviors on Social Media}\label{sec:model-behaviors}
%\paragraph{Post.}
At each access, agent $i$ creates a post $p_i$ reflecting its current
opinion $O_i(t)$ with probability $\post_i(t)$. This post is made
available to every current follower of $i$ and is delivered via their
feeds, as described below.
%\paragraph{Feed.}
Agent $i$'s {\em feed} $\feed_i(t)$ is assembled
in a strictly reverse-chronological order from the agents that $i$ follows,
comprising the $\maxfeed$ most recent original or reposted posts.
%\paragraph{Repost.} 
Upon reading its feed, agent $i$ reposts each post it finds tolerable,
i.e., within its bounded confidence threshold
$\BoundedConfidence_i(t)$, independently with probability $P_r$. A
post that falls outside $\BoundedConfidence_i(t)$ may still be
reposted, but only at the much lower rate $P_r^{\text{out}}\ll P_r$,
reflecting the possibility of occasional out-of-tolerance
amplification, such as quoting a disagreeable post to react to it.
% A repost retains the original post's opinion 
% content and original authorship (i.e. agent index), 
% and is delivered to every follower of the reposting 
% agent exactly as an original post would be.
\par


%\paragraph{Follow.} 
After reading its feed, agent $i$ collects, as candidates, the authors
of all posts within its bounded confidence threshold whom it does not
already follow and has not blocked. With probability $P_f$, it follows
one candidate chosen uniformly at random from these candidates. The
maximum number of agents that any agent can follow is fixed at
$\maxfollower$ throughout the experiments, based on the concept of
\emph{limited attention}: although a user may nominally follow
hundreds of accounts, the number of information sources that can
effectively influence it at any given time is constrained by human
cognitive limits and information
overload~\cite{goncalves2011,weng2012}. Once at capacity, a new follow
immediately evicts the currently followed agent whose opinion is
farthest from $i$'s own.
\par


%\paragraph{Unfollow and block.} 
At each access, agent $i$ first makes a single probabilistic decision
with probability $P_u$ about whether it will act on any
out-of-tolerance content encountered during that access. If this
decision succeeds, $i$ scans its feed for the first post whose author
it follows but whose opinion lies outside $\BoundedConfidence_i(t)$;
if such a post exists, $i$ unfollows and blocks that
author. Otherwise, if the feed contains any out-of-tolerance post from
an author that $i$ does not follow, $i$ blocks one such author, chosen
uniformly at random. Therefore, at most one follow edge is severed, or
one author blocked, per access. Thus, $P_u$ sets the population's
overall propensity to prune out-of-tolerance ties. We treat it as the
model's central structural knob, governing how readily the network
fragments into mutually invisible clusters as opposed to remaining
broadly connected.
\par

Posting probability $\post_i(t)$ is deliberately kept independent of
opinion extremity $\vert O_i(t)\vert$ throughout the rules
above. This independence ensures that posting and reposting behaviors
serve directly as the primary outcomes of interest in this study: it
is this surfaced, vocalized activity that other users actually
perceive and are influenced by, rather than the unexpressed internal
opinions themselves.
\par


\begin{table}[t]\Blue
\centering
\caption{Model parameters and their default values.}\label{tab:param}
\tabcolsep=5pt
\begin{tabular}{lll}
\toprule
Parameter & Symbol & Default value \\
\midrule
Probability of accessing the platform & $P_a$ & 0.1 \\
Probability of following & $P_f$ & 0.01 \\
Probability of reposting (in/out of tolerance) & $P_r$, $P_r^{\text{repost}}$ & 0.1, 0.01 \\
Probability of unfollowing/blocking & $P_u$ & 0.05 (variable) \\
Maximum number of followees & $\maxfollower$ & 43.9 \\
Feed capacity & $\maxfeed$ & 15 \\
Confidence decline/recovery rate & $\delta_{\BoundedConfidence}$, $\gamma_{\BoundedConfidence}$ & 0.5, $10^{-4}$ \\
Bounded confidence floor/ceiling & $\BoundedConfidence_{\mathrm{floor}},\BoundedConfidence_{\mathrm{ceil}}$ & 0.2, 1.0 \\
Vocal comfort radius & $\VCR$ & 0.1 (variable) \\
Posting-probability relaxation rate & $\eta$ & 0.1 \\
Posting-probability floor/ceiling & $\post_{\min},\post_{\max}$ & 0.01, 1.0 \\
Stubbornness & $\Stubbornness$ & 0.82 \\
\bottomrule
\end{tabular}\Black
\end{table}


\subsection{Modeling the Moderate Hub Intervention}\label{sec:model-intervention}
The manipulation central to this study exogenously fixes the opinions
of a designated set of highly reachable agents. At $t=0$, we rank all
agents by their initial follower count and select the
$N_{\mathrm{hub}}=60$ highest-ranked as the hub pool. Follower count
is chosen here because it is currently the only channel through which
reach operates in this model. The number $N_{\mathrm{hub}}=60$ is
fixed well above the several dozen agents to which the network
generator's heavy tail (Section~\ref{sec:experiments}) gives
comparably large in-degree, so that the hub pool conservatively covers
this entire plateau of high-reach agents. This pool is partitioned
into five categories, symmetric about the moderate center and weighted
toward it: ten agents each at the extreme positions $O=-1.0$ and
$O=+1.0$; ten agents each at the near-extreme positions $O=-0.4$ and
$O=+0.4$; and twenty agents at the moderate position $O=0.0$, the
moderate hubs whose activity is the object of this study --- with
these five target values being representative of the \emph{extreme},
\emph{near-extreme}, and \emph{moderate} camps defined in
Section~\ref{sec:model-abm}, respectively. Every hub agent's intrinsic
and current opinion is reset once, at $t=0$, to its assigned target
and held fixed thereafter by exempting it from
Eq.~\eqref{eq:social-influence}; its stubbornness, posting
probability, and bounded confidence dynamics are otherwise left
unmodified.
\par

The model further supports silencing a hub category. A silenced hub is
prevented from posting original content or reposting, while continuing
to read its feed and to follow or unfollow exactly as before, so its
follow relationships evolve identically to those under an active
hub. Contrasting an active-hub condition against an otherwise
identical hub-silenced condition therefore isolates the effect
specifically attributable to the hub's own vocal activity. The
specific experimental arms constructed from this mechanism are
described in Section~\ref{sec:experiments}.
\par

%
\Black%%{}######
%

\section{Experimental Design}\label{sec:experiments}

The initial follow graph is generated by a directed
Holme--Kim growing network process~\cite{holme2002} 
with $m=15$, $A_0=4$ 
and $p_t=0.05$, which produces 
synthetic heavy-tailed network mimicing real-world 
asymmetric web social networks.
Most of the parameter values listed in
Table~\ref{tab:param}, together with the
network-generator parameters, were 
fixed by a joint parameter
sweep against several empirical targets drawn from
prior studies of Twitter/X
highliging heavy-tailed degree distribution~\cite{kwak2010}, 
high modularity of the reposting network~\cite{conover2011}, 
posting dynamics weighted by opinion strength~\cite{buntain2025}, 
thereby enable the model to reproduce realistic polarized posting dynamics
on social media. 
The two
parameters swept as experimental factors below,
$P_u$ and $\VCR$, are excluded from this
calibration and instead varied directly.

Within this calibrated network, 
we manipulate the moderate hub category
introduced in Section~\ref{sec:model-intervention}
across two arms: a \emph{control} arm, in which
all five hub categories remain active throughout,
and an \emph{M-silent} arm, in which only the
moderate hub category is silenced while the
extreme and near-extreme categories remain active
in both arms. Because silencing leaves every
hub's follow relationships unchanged
(Section~\ref{sec:model-intervention}), the isolated
difference between the two arms is whether the
moderate hubs actually post and repost, so that
any difference in outcome between them can be
attributed to the moderate hubs' own vocal
activity.

We cross this arm manipulation with a full
factorial sweep of the two parameters: the
vocal comfort radius $\VCR\in\{0.1,0.15,0.2\}$,
capped at the bounded confidence floor
$\BoundedConfidence_{\mathrm{floor}}=0.2$
(Table~\ref{tab:param}) reflecting the notion that 
an agent's
tolerance can never decline below what it needs
to feel comfortable enough to speak, and
the unfollow/block propensity $P_u\in\{0.05,0.15\}$,
giving six parameter cells crossed with the two
arms above. Each of these $6\times2=12$ conditions
is simulated with $N=1000$ agents over $40{,}000$
steps, replicated across $n=15$ independent
network seeds per condition.

The outcome we compare between arms is the
posting activity of each 5 camp
(Section~\ref{sec:model-abm}), restricted to
non-hub agents, since the hub agents' own opinions
and activity are exogenously fixed rather than a
response to be measured.
Reposts are excluded from this
count, so that the outcome reflects agents' own
willingness to voice an opinion rather than their
willingness to amplify someone else's.

\section{Results}\label{sec:results}

\begin{figure}[t]
\centering
\fbox{\parbox[c][6cm]{0.9\linewidth}{\centering
TODO: heatmap figure --- moderate / near-extreme /
extreme camp posting-rate difference (M-silent
minus control), one panel per camp, $P_u$ on one
axis and $\VCR$ on the other}}
\caption{Per-camp posting-rate difference
(M-silent minus control) across the $P_u\times\VCR$
grid, for the moderate, near-extreme, and extreme
camps.}
\label{heatmap}
\end{figure}

Fig.~\ref{heatmap} shows heatmaps of the difference
in per-step per-capita posting counts for each of
the three opinion camps, averaged over steps
[39000, 40000], comparing \emph{M-silent} against
\emph{control}.
Steadily, the moderate camp boosts its posting
activity when it can perceive aligned-opinion
hubs' vocality, whereas the near-extreme camp
shrinks its posting by contrast.
Compared to these fairly intuitive results, the
extreme camp does not show a linear result: its
posting activity is reduced (a negative value in
the heatmap) when $\VCR$ is relatively low, for
each $P_u$ row, and it crosses the sign of its
effect at a different value of $\VCR$ depending on
the $P_u$ value.
Equivalently, reversing the contrast from
\emph{M-silent -- control} to \emph{control --
M-silent} lets us easily interpret the same difference
from the complementary perspective of the moderate
hub's activity rather than its absence. 
From this point of view, focusing on the presence
of moderate hubs, under the same propensity to
prune out-of-tolerance ties ($P_u$), moderate hubs
facilitate the extreme camp's posting activity
when agents speak only in the presence of
sufficient nearby agreement (low $\VCR$), and
mediate it in the opposite environment.

Also, $P_u$ can be interpreted as a modulator that
shifts the $\VCR$ cross-sign point, the watershed
value at which moderate hubs' vocality switches
between acting as a facilitator or a mediator for
the extreme camp.
In the extreme camp result of Fig.~\ref{heatmap},
the effect size in higher $P_u$ rows stays
negative across a wider range of low $\VCR$ before
turning positive at higher $\VCR$, compared with
lower $P_u$ rows.
This confirms that $P_u$ governs where this
regime shift occurs. 

\section{Discussion}\label{sec:discussion}
Our intuitive results for moderate and near-extreme camps ---
the former amplified and the latter silenced under moderate hubs 
vocalization--- 
show its linearness seemingly because the both camp basically 
stay inside the reach of moderate hubs. As an agent in our 
model modulate their posting probability only depending on 
the extent to which it feel comfortable in its feed, 
peripheral moderate agents receive more posts from moderate 
hubs in \emph{control} arm which lead to their higher 
comfortable feed resulting in larger vocalization. 
Subsequently, near-extreme camp much frequently encounter 
moderate-opinion posts when moderate hubs activate, thereby 
their feeds relatively reduce its critically aligned near-extreme 
posts which suffcietnly amplify their vocality. 
This silence effect for near-extreme could happen because 
posts from moderate hubs is mechanically able to reach their 
feeds through the networks. 

The regime shift results of the effect for extreme camp 
seems to be organized above a single mechanistic hypothesis: 
moderate hubs act as a mediator when extreme camp can 
be considered to homophilic to near-extreme camp, and 
act as a facilitator when considered to heterophic to 
near-extreme camp. 
First, the watershed depending on two variables appear to 
indicate this homophily-or-heterophy logic between 
near-extremes and extremes. 

Let's see results of the column where $\VCR$ is between 
0.1 and 0.2. This means that the role of moderate hubs 
whether as a facilitator or a mediator could change 
in the same $\VCR$ depending on the disconnecting 
propensity of the population($P_u$). 
In order to organize this crossing-point, 
we focus on the distinct effect of $P_u$ and $\VCR$, 
which govern structural and perceptional homophilic range. 
First, for $P_u$, in its higher value, social network 
prone to be segregated, thereby the range to which extreme 
camp can reach and perceive homophilic 
in neighbor near-extreme camp 
shall be structurally limited. 
Upon this structural layer, for $\VCR$, 
the range of homophily reach from extremes 
toward near-extremes is to be larger 
when $P_u$ is low because of its 
much many absolute number of connections with them. 

Then, the regime shift in the extremes could be 
explained by organizing this two layer 
homophily hypothesis and 
the moderate hubs' silence effect over near-extremes discussed above. 
As near-extremes are always silenced under moderate hubs' vocalization, 
extremes seem to be also silenced when these two camps 
considered to be homophilic governed by two layer 
mechanism. 
On the contrary, when these two camps are not considered 
to be sufficiently homophilic, 
the vocality of extrems is amplified under the silence of 
near-extreme that let extremes be more isolated and purified. 
More mechanically, because the effect of moderate hubs 
vocalization here 
in this model can be considered to replace some amount of 
posts in agents' feeds with their moderate-opinion posts
whether through posting or reposting, 
how far in a social-network-wise the hubs can reach and influence 
decides their replacement reachability. 
Interestingly, we find an analogous sign inversion related to 
moderate hubs effect in a former research~\cite{betts2023}, 
which indicates 
that moderate influencer could worsen polarization 
depending on the influencer's own assimilation tolerance, 
independently converging on the same qualitative phenomenon 
regarding moderate hubs. 

\section{Conclusion}\label{sec:conclusion}
Although the separation by fifth bins with agents' opinions 
degenerate this continuous ripple-ish effect, this formalization 
sufficiently gives us the implication of counter-intuitive 
effect by moderate hubs' vocality. -> as a limitation and 
hint for further reserach.

\bibliographystyle{spmpsci}
\bibliography{refs}
\end{document}
